\chapter{วิธีการดำเนินการวิจัย}
\label{chapter:experiment}

ขั้นตอนการศึกษาและพัฒนาระบบผู้ช่วยอัจฉริยะเพื่อการศึกษา (UniAssist AI) แบ่งออกเป็นขั้นตอนหลักดังต่อไปนี้

\begin{enumerate}
    \item การเตรียมข้อมูลและการสร้างฐานข้อมูล (Data Preparation and Database Construction)
    \item การประเมินและคัดเลือกโมเดลภาษาสำหรับการสกัดข้อมูล (Model Evaluation and Selection)
    \item การพัฒนาแกน Text-to-SQL Chatbot
    \item การพัฒนาระบบประมวลผลฟังก์ชันพิเศษ (คำนวณเกรด ประเมินความเสี่ยง ทุนการศึกษา)
    \item การพัฒนาส่วนติดต่อผู้ใช้ ระบบแสดงผล และการยืนยันตัวตน
\end{enumerate}

\section{ภาพรวมสถาปัตยกรรมของระบบ}

ระบบ UniAssist AI ถูกพัฒนาในรูปแบบเว็บแอปพลิเคชันด้วยเฟรมเวิร์ก Flask (ภาษา Python) โดยแบ่งการทำงานออกเป็น 3 ชั้นหลัก ได้แก่

\begin{enumerate}
    \item \textbf{ชั้นข้อมูล (Data Layer)} ฐานข้อมูลเชิงสัมพันธ์ SQLite ที่จัดเก็บข้อมูลหลักสูตร รายวิชา ระเบียบข้อบังคับ ค่าเกรด ทุนการศึกษา และข้อมูลนักศึกษา โดยระบบเข้าถึงฐานข้อมูลนี้แบบอ่านอย่างเดียว (read-only) ในทุกส่วนงาน
    \item \textbf{ชั้นประมวลผล (Application Layer)} ประกอบด้วยแกน Text-to-SQL (\texttt{text\_to\_sql\_chatbot.py}) ตรรกะการคำนวณเกรดและประเมินความเสี่ยง (\texttt{assist\_logic.py}) และการยืนยันตัวตน (\texttt{google\_auth.py}) โดยเรียกใช้บริการโมเดลภาษาขนาดใหญ่ผ่าน API ที่เข้ากันได้กับมาตรฐาน OpenAI
    \item \textbf{ชั้นแสดงผล (Presentation Layer)} เว็บอินเทอร์เฟซสำหรับนักศึกษาและแดชบอร์ดสำหรับอาจารย์ที่ปรึกษา
\end{enumerate}

การออกแบบให้ทุกส่วนเข้าถึงฐานข้อมูลแบบอ่านอย่างเดียว และให้โมเดลภาษาเป็นเพียงผู้ ``สร้างคำสั่งค้นข้อมูล'' ไม่ใช่ผู้เก็บข้อเท็จจริงไว้ในตัวเอง เป็นหลักการสำคัญที่ทำให้คำตอบของระบบสามารถตรวจสอบย้อนกลับไปยังข้อมูลจริงในฐานข้อมูลได้เสมอ

\section{การเตรียมข้อมูลและการสร้างฐานข้อมูล}

\subsection{การวิเคราะห์ความต้องการ}

ในขั้นตอนแรกได้ทำการวิเคราะห์ปัญหาและความต้องการของกลุ่มเป้าหมาย คือ นักศึกษาและอาจารย์ที่ปรึกษา เพื่อกำหนดขอบเขตฟังก์ชันหลัก เช่น การเข้าถึงข้อมูลหลักสูตร การวางแผนการเรียน การคำนวณเกรด และการติดตามความเสี่ยงทางการเรียน

\subsection{การรวบรวมข้อมูล}

ข้อมูลที่ใช้ในโครงงานเป็นข้อมูลทุติยภูมิที่รวบรวมจากแหล่งข้อมูลทางการของสถาบัน ได้แก่

\begin{enumerate}
    \item \textbf{เอกสารรายละเอียดหลักสูตร (มคอ.2)} ของทั้ง 4 หลักสูตรในระดับปริญญาตรีของคณะเทคโนโลยีสารสนเทศ ได้แก่ IT, DSBA, BIT และ AIT ซึ่งระบุโครงสร้างหลักสูตร รายวิชา หน่วยกิต แผนการเรียน และเงื่อนไขบังคับก่อน
    \item \textbf{ตารางสอนและข้อมูลจากสำนักทะเบียน} ผ่านการเรียกใช้บริการ (API) ของระบบทะเบียน โดยข้อมูลปี/ภาคการศึกษา/คณะ/หลักสูตร ถูกส่งเป็นพารามิเตอร์ของคำขอ
    \item \textbf{ระเบียบและข้อบังคับการศึกษา} เช่น เกณฑ์การพ้นสภาพ ภาคทัณฑ์ เกียรตินิยม และเงื่อนไขการลงทะเบียน
\end{enumerate}

\subsection{การสกัดข้อมูลด้วยโมเดลภาษา}

เนื่องจากข้อมูลหลักสูตรในเอกสาร มคอ.2 อยู่ในรูปแบบเอกสารกึ่งโครงสร้าง (ตารางและข้อความบรรยาย) ที่ยากต่อการแปลงเป็นฐานข้อมูลโดยตรง จึงได้ประยุกต์ใช้โมเดลภาษาขนาดใหญ่ในการสกัดข้อมูล (Information Extraction) ให้อยู่ในรูปแบบ JSON ที่มีโครงสร้างชัดเจน กระบวนการสกัดถูกทดลองกับโมเดลหลายตัว ได้แก่ Qwen, Gemini, DeepSeek และ Typhoon เพื่อเปรียบเทียบคุณภาพ จากนั้นนำผลลัพธ์ JSON ที่ได้มาประมวลผลและนำเข้าสู่ฐานข้อมูลด้วยชุดสคริปต์ที่พัฒนาขึ้น (\texttt{build\_qwen\_db.py}, \texttt{build\_chatbot\_db.py} และ \texttt{build\_extra\_tables.py}) โดยสคริปต์เหล่านี้ออกแบบให้ทำงานซ้ำได้ (idempotent) และแยกส่วนของข้อมูลที่ผ่านการตรวจสอบด้วยมือ (Ground Truth) ออกจากส่วนที่สกัดด้วยโมเดล

\subsection{การออกแบบฐานข้อมูล}

ฐานข้อมูลถูกออกแบบในรูปแบบเชิงสัมพันธ์ ประกอบด้วยตารางหลักที่เชื่อมโยงกันด้วยรหัสหลักสูตร (\texttt{program\_id}) และรหัสวิชา ตารางที่ \ref{tab:schema} แสดงตารางสำคัญในฐานข้อมูลและหน้าที่ของแต่ละตาราง

\begin{table}[ht]
    \centering
    \caption{ตารางสำคัญในฐานข้อมูลของระบบ UniAssist AI}
    \label{tab:schema}
    \begin{tabularx}{\linewidth}{l l X}
        \toprule
        \textbf{ตาราง} & \textbf{จำนวนแถว} & \textbf{หน้าที่} \\
        \midrule
        \texttt{programs} & 4 & ข้อมูลหลักสูตร (ชื่อ ปริญญา หน่วยกิตรวมที่ต้องเรียนจบ) \\
        \texttt{courses} & 300 & รายวิชาทั้งหมดในหลักสูตร พร้อมหน่วยกิตและประเภท \\
        \texttt{academic\_plan\_courses} & 583 & แผนการเรียนรายเทอม (แยกแผนปกติ/สหกิจ และสายวิชา) \\
        \texttt{course\_prerequisites} & 30 & เงื่อนไขรายวิชาบังคับก่อน (Prerequisite) \\
        \texttt{curriculum\_structure} & 16 & โครงสร้างหมวดวิชาและจำนวนหน่วยกิตของแต่ละหมวด \\
        \texttt{rules} & 108 & ระเบียบ/เกณฑ์ (รีไทร์ ภาคทัณฑ์ เกียรตินิยม การลงทะเบียน ฯลฯ) \\
        \texttt{grade\_scale} & 12 & ค่าระดับคะแนน (A=4.0, B+=3.5, \ldots) และสถานะการนับ GPA \\
        \texttt{scholarships} & 11 & ข้อมูลทุนการศึกษาและเงื่อนไขคุณสมบัติ \\
        \texttt{students} & 17 & ข้อมูลนักศึกษา (หลักสูตร อาจารย์ที่ปรึกษา GPAX สถานะ) \\
        \texttt{student\_term\_gpa} & 28 & ผลการเรียนเฉลี่ยรายภาคการศึกษาของนักศึกษา \\
        \bottomrule
    \end{tabularx}
\end{table}

\section{การประเมินและคัดเลือกโมเดลภาษาสำหรับการสกัดข้อมูล}

เนื่องจากคุณภาพของฐานข้อมูลขึ้นอยู่กับความถูกต้องของการสกัดข้อมูลโดยตรง จึงได้ออกแบบกระบวนการประเมินเพื่อคัดเลือกโมเดลที่เหมาะสมที่สุด โดยเปรียบเทียบผลการสกัดของแต่ละโมเดลกับเอกสาร มคอ.2 ฉบับจริง (ใช้เป็น Ground Truth) ผ่านตัวชี้วัด 2 มิติ ได้แก่

\begin{enumerate}
    \item \textbf{ความครบถ้วน (Coverage)} สัดส่วนของรายวิชาที่โมเดลสามารถดึงมาได้ครบเมื่อเทียบกับจำนวนรายวิชาทั้งหมดในเอกสาร
    \item \textbf{ความถูกต้องของแอตทริบิวต์ (Attribute Accuracy / Agreement)} เมื่อดึงรายวิชามาได้แล้ว ค่าของแต่ละแอตทริบิวต์ตรงกับเอกสารจริงหรือไม่ โดยพิจารณา 7 แอตทริบิวต์ ได้แก่ ชื่อวิชา (ภาษาไทย), ชื่อวิชา (ภาษาอังกฤษ), หน่วยกิต, ปี/ภาคการศึกษา, หมวดวิชา, ประเภทวิชา และรหัสวิชา
\end{enumerate}

การประเมินดำเนินการครอบคลุมทั้ง 4 หลักสูตร โดยผลการประเมินโดยละเอียดและการคัดเลือกโมเดลจะนำเสนอในบทที่ \ref{chapter:result}

\section{การพัฒนาแกน Text-to-SQL Chatbot}

แกนหลักของระบบทำหน้าที่แปลงคำถามภาษาไทยของผู้ใช้ให้เป็นคำสั่ง SQL และเรียบเรียงผลลัพธ์กลับเป็นคำตอบภาษาไทย โดยมีลำดับการทำงาน (Pipeline) ดังนี้

\begin{enumerate}
    \item \textbf{การสร้างบริบทของ Schema} ระบบอ่านโครงสร้างฐานข้อมูลจริง (Schema Introspection) มาสร้างเป็นคำอธิบายประกอบในคำสั่งระบบ (System Prompt) พร้อมด้วยหมายเหตุสำคัญเกี่ยวกับข้อมูล เช่น หน่วยกิตจบของแต่ละหลักสูตร กฎการเลือกแผนการเรียน และข้อควรระวังในการเขียน SQL การอ่าน schema จากฐานข้อมูลจริงช่วยป้องกันปัญหาคำสั่งระบบล้าสมัยเมื่อโครงสร้างข้อมูลเปลี่ยนแปลง
    \item \textbf{การจำแนกประเภทคำถาม (Intent Classification)} ระบบให้โมเดลตัดสินว่าคำถามเป็นคำถามที่ต้องดึงข้อมูล (สร้าง SQL) หรือเป็นการสนทนาทั่วไป (CHAT) หากไม่สามารถตอบจากข้อมูลได้ โมเดลจะคืนค่า \texttt{NO\_ANSWER} แทนการเดา
    \item \textbf{การสร้างคำสั่ง SQL (SQL Generation)} เมื่อเป็นคำถามดึงข้อมูล โมเดลจะสร้างคำสั่ง SQL โดยตั้งค่า temperature เท่ากับ 0 เพื่อให้ผลลัพธ์มีความคงเส้นคงวา
    \item \textbf{การตรวจสอบความปลอดภัยของคำสั่ง (SQL Sanitization)} ก่อนรันทุกคำสั่ง ระบบจะตรวจสอบว่า (ก) เป็นคำสั่ง \texttt{SELECT} หรือ \texttt{WITH} เท่านั้น (ข) ไม่มีคำสั่งเขียน/แก้ไขข้อมูล (ค) ป้องกันการรันหลายคำสั่งพร้อมกัน (Multiple Statements) และ (ง) เติม \texttt{LIMIT} อัตโนมัติหากไม่มี เพื่อป้องกันการดึงข้อมูลทั้งตาราง
    \item \textbf{การประมวลผลคำสั่ง (Query Execution)} รันคำสั่งบนการเชื่อมต่อฐานข้อมูลแบบอ่านอย่างเดียว (\texttt{file:...?mode=ro}) ทำให้แม้โมเดลจะสร้างคำสั่งที่เป็นอันตราย ระบบก็ไม่สามารถแก้ไขข้อมูลได้โดยเด็ดขาด
    \item \textbf{การเรียบเรียงคำตอบ (Answer Phrasing)} นำผลลัพธ์ที่ได้ส่งกลับให้โมเดลเรียบเรียงเป็นคำตอบภาษาไทยที่กระชับ โดยกำชับให้อ้างอิงจากผลลัพธ์เท่านั้น หากไม่มีข้อมูลให้ตอบตรง ๆ ว่าไม่พบ
    \item \textbf{กลไกสำรอง (Fallback)} หากรอบแรกได้ \texttt{NO\_ANSWER} หรือคำสั่งผิดพลาด ระบบจะลองสร้าง SQL ใหม่อีกครั้ง (retry) และหากยังไม่สำเร็จจะสลับไปโหมดสนทนาทั่วไป ซึ่งตอบโดยยึด (Grounding) จากสรุปข้อเท็จจริงในฐานข้อมูล เพื่อลดการตอบนอกบริบทและการสร้างข้อมูลเท็จ
\end{enumerate}

การออกแบบให้ระบบ ``แสดงคำสั่ง SQL ที่โมเดลเขียนเสมอ'' มีวัตถุประสงค์เพื่อให้สามารถตรวจสอบได้ว่าคำตอบมาจากการค้นข้อมูลใด และเพื่อใช้เปรียบเทียบกับผลเฉลย (Ground Truth) ในการประเมินความถูกต้องของระบบ

\section{การพัฒนาระบบคำนวณเกรดและประเมินความเสี่ยง}

ระบบคำนวณเกรดถูกออกแบบให้อ้างอิงค่าจากฐานข้อมูลเป็นแหล่งข้อมูลหลักเดียว (Single Source of Truth) เพื่อหลีกเลี่ยงการกำหนดค่าคงที่ (Hardcode) ที่อาจล้าสมัย โดยมีขั้นตอนดังนี้

\begin{enumerate}
    \item \textbf{การแยกวิเคราะห์ผลการเรียน (Transcript Parsing)} ผู้ใช้สามารถวางข้อความผลการเรียน (Transcript) หรือแก้ไขรายวิชาในตาราง ระบบจะแยกวิเคราะห์รหัสวิชา หน่วยกิต และเกรดออกมา โดยรองรับรูปแบบหัวข้อภาคการศึกษาที่หลากหลาย
    \item \textbf{การคำนวณ GPA} คำนวณคะแนนเฉลี่ยรายภาคการศึกษาและคะแนนเฉลี่ยสะสม โดยใช้ค่าระดับคะแนนจากตาราง \texttt{grade\_scale} และตัดรายวิชาที่ไม่นับหน่วยกิต GPA (เช่น S, U, T, I และวิชาที่ถอน W) ออกจากการคำนวณ
    \item \textbf{การจัดสถานะและประเมินความเสี่ยง} เปรียบเทียบ GPAX กับเกณฑ์ในตาราง \texttt{rules} (เช่น เกณฑ์รีไทร์ ภาคทัณฑ์ และเกียรตินิยม) เพื่อจัดสถานะและประเมินความเสี่ยง จากนั้นใช้ LLM สรุปสถานะและให้คำแนะนำเชิงปฏิบัติ เนื่องจากข้อความกฎในฐานข้อมูลบางส่วนจัดเก็บตัวเลขเป็นเลขไทย ระบบจึงมีการแปลงเลขไทยเป็นเลขอารบิกก่อนนำไปเปรียบเทียบ
\end{enumerate}

\section{การพัฒนาระบบทุนการศึกษา}

ระบบข้อมูลทุนการศึกษาดึงข้อมูลจากตาราง \texttt{scholarships} ซึ่งประกอบด้วยชื่อทุน ผู้ให้ทุน จำนวนเงิน เกณฑ์ GPA คุณสมบัติ และกำหนดเวลารับสมัคร โดยระบบสามารถใช้ LLM ช่วยวิเคราะห์คุณสมบัติเบื้องต้นของนักศึกษาเทียบกับเงื่อนไขของทุนแต่ละประเภท และสรุปทุนที่นักศึกษามีสิทธิ์สมัครได้

\section{การพัฒนาแดชบอร์ดสำหรับอาจารย์ที่ปรึกษา}

แดชบอร์ดสำหรับอาจารย์ที่ปรึกษาแสดงภาพรวมของนักศึกษาในความดูแล โดยดึงข้อมูลจากตาราง \texttt{students} และ \texttt{student\_term\_gpa} มาสรุปผลการเรียนเฉลี่ยรายภาคการศึกษา และจำแนกกลุ่มความเสี่ยงของนักศึกษา เพื่อให้อาจารย์สามารถระบุนักศึกษาที่ควรได้รับการดูแลเป็นพิเศษได้อย่างรวดเร็ว

\section{การยืนยันตัวตนและความปลอดภัย}

เนื่องจากนักศึกษาเป็นผู้พิมพ์คำถามเข้าสู่ระบบโดยตรง การรักษาความปลอดภัยจึงเป็นประเด็นสำคัญ ระบบจึงออกแบบการป้องกันหลายชั้น ได้แก่

\begin{enumerate}
    \item \textbf{การยืนยันตัวตนด้วย Google OAuth 2.0} ผู้ใช้ต้องเข้าสู่ระบบด้วยบัญชี Google ที่จำกัดเฉพาะโดเมนของสถาบัน และแบ่งสิทธิ์การใช้งานตามบทบาท (นักศึกษา / อาจารย์ที่ปรึกษา / ผู้พัฒนา) โดยข้อมูลผู้ใช้ถูกจัดเก็บใน session cookie ที่มีการลงลายเซ็น
    \item \textbf{การจำกัดคำสั่งฐานข้อมูล} เชื่อมต่อฐานข้อมูลแบบอ่านอย่างเดียว ตรวจสอบให้เป็นคำสั่ง \texttt{SELECT}/\texttt{WITH} เดี่ยว บล็อกคำสั่งเขียนข้อมูล ป้องกันการรันหลายคำสั่ง และเติม \texttt{LIMIT} อัตโนมัติ ตามที่กล่าวในหัวข้อแกน Text-to-SQL
    \item \textbf{การป้องกันข้อมูลลับ} ค่าคีย์ลับต่าง ๆ (เช่น Google Client Secret และ API Key ของบริการ LLM) ถูกอ่านและใช้งานเฉพาะฝั่งเซิร์ฟเวอร์เท่านั้น
\end{enumerate}
